% ---------------------------------------------------------------------
% Appendix B
% ---------------------------------------------------------------------

\begin{cnread}
本附录把「曲率形式」与「和乐密度」之间的因子 $2$ 老老实实算一遍。
结论是：交叉比和乐 $\Lambda(\square)$ 的一阶展开密度是
$-(1+ab)^{-2}\,\mathrm da\wedge\mathrm db$，而 Paper~IV 备注 4.12 给出的
约化联络 $\omega_H$ 的曲率是 $+2(1+ab)^{-2}\,\mathrm da\wedge\mathrm db$。
两者都对，但\textbf{规范化不同}：传递元是那条联络"焊接项取一半"后的和乐。
任何把辐射功率等同到曲率积分的说法，必须先把这个因子定下来。
\end{cnread}

\subsection{The curvature of the reductive connection}

Let $f=(1+ab)^{-1}$ and consider the non-closed part of the connection
\eqref{eq:connection},
\begin{equation}\label{eq:omega2}
  \omega_2=f\bigl(a\,\mathrm db-b\,\mathrm da\bigr).
\end{equation}
Because $\mathrm d(a\,\mathrm db-b\,\mathrm da)=2\,\mathrm da\wedge\mathrm db$,
\begin{equation}\label{eq:domega}
  \mathrm d\omega_2
  =\mathrm df\wedge\bigl(a\,\mathrm db-b\,\mathrm da\bigr)+2f\,\mathrm da\wedge\mathrm db .
\end{equation}
Now $\mathrm df=-(1+ab)^{-2}(b\,\mathrm da+a\,\mathrm db)$, so
\begin{align}
  \mathrm df\wedge\bigl(a\,\mathrm db-b\,\mathrm da\bigr)
  &=-(1+ab)^{-2}\bigl[(b\,\mathrm da+a\,\mathrm db)\wedge(a\,\mathrm db-b\,\mathrm da)\bigr]
  \notag\\
  &=-(1+ab)^{-2}\cdot 2ab\,\mathrm da\wedge\mathrm db ,
  \label{eq:df-term}
\end{align}
since $(b\,\mathrm da+a\,\mathrm db)\wedge(a\,\mathrm db-b\,\mathrm da)
=ab\,\mathrm da\wedge\mathrm db+ab\,\mathrm da\wedge\mathrm db$.  Combining
\eqref{eq:domega} and \eqref{eq:df-term},
\begin{equation}\label{eq:curvature}
  \mathrm d\omega_H=\mathrm d\omega_2
  =2\,\mathrm da\wedge\mathrm db\Bigl(\frac1{1+ab}-\frac{ab}{(1+ab)^{2}}\Bigr)
  =\frac{2\,\mathrm da\wedge\mathrm db}{(1+ab)^{2}} .
\end{equation}
Check~V4.1 verifies the coefficient
$\partial_a(fa)+\partial_b(fb)=2(1+ab)^{-2}$ symbolically.

\subsection{The density of the transport}

Independently, the first-order expansion of the cross-ratio holonomy around a
coordinate rectangle of side lengths $\varepsilon$ and $\eta$ based at $(a_0,b_0)$
is, by \eqref{eq:cross-ratio},
\begin{equation}\label{eq:log-lambda}
  \log\Lambda(\square)
  =-\frac{\varepsilon\eta}{(1+a_0b_0)^{2}}+O(3)
  =-\iint_{\square}\frac{\mathrm da\wedge\mathrm db}{(1+ab)^{2}}+O(3),
\end{equation}
which is Check~V3.1.  Comparing \eqref{eq:curvature} with \eqref{eq:log-lambda},
the transport density is minus one half of the connection's curvature density.

\subsection{Reconciliation}

Three statements are exact and should be kept apart.

\begin{enumerate}[label=(\roman*)]
  \item The cross-ratio formula \eqref{eq:cross-ratio} is exact and is the
  holonomy of the idempotent transport of Definition~\ref{def:torsion-transport}
  (Checks~V1.1--V2.5).
  \item The connection \eqref{eq:connection} has curvature
  $2(1+ab)^{-2}\mathrm da\wedge\mathrm db$ (Check~V4.1).
  \item The first-order log-holonomy density is
  $-(1+ab)^{-2}\mathrm da\wedge\mathrm db$ (Check~V3.1).
\end{enumerate}

The holonomy of a connection with curvature $\Omega$ around a small positively
oriented loop is $\exp\bigl(-\int_\Sigma\Omega\bigr)$ in the abelian
approximation.  Statement (ii) would therefore give
$\log\Lambda=-\int_\Sigma 2(1+ab)^{-2}\mathrm da\wedge\mathrm db$, which differs
from (iii) by the factor $2$.  The resolution is that the idempotent transport is
the holonomy of the connection whose reductive second term is
$\tfrac12 f(a\,\mathrm db-b\,\mathrm da)$, i.e. of
$\omega_H^{(1/2)}:=\Phi^{-1}\mathrm d\Phi+\tfrac12 f(a\,\mathrm db-b\,\mathrm da)$.
Its curvature is $\tfrac12\cdot2(1+ab)^{-2}\mathrm da\wedge\mathrm db$, whose
negative integral is exactly \eqref{eq:log-lambda}.

\begin{remark}[Which normalisation is \lq\lq right\rq\rq]
Neither is more correct.  The factor is the difference between a connection
normalised so that its transverse term is the full $(1+ab)^{-1}(a\,\mathrm
db-b\,\mathrm da)$ and one normalised so that the transport's multiplicative
component is the raw cross-ratio.  What matters is that a physical
identification of $\Lambda(\square)-1$ with a radiated power must state which
convention it uses; otherwise the same history gives two answers differing by a
factor two.

The two exact statements are now in Paper~IV itself: Remark~4.12 carries the
curvature identity $\mathrm d\omega_H=2(1+ab)^{-2}\mathrm da\wedge\mathrm db$
together with the correct Borel reduction of the solder component
(equation~\eqref{eq:solder-borel} above), and Proposition~4.17 identifies the
connection whose holonomy reproduces $\Lambda(\square)$ to first order.  The
choice between the two normalisations is left open there, as it is here.
\end{remark}

\begin{remark}[Sign]
\label{rem:sign}
The two densities also differ in sign, because the loop is traversed in the
orientation that makes $\mathrm da\wedge\mathrm db$ positive while the
log-holonomy is negative.  Both statements are checked, and the signs are
consistent with the holonomy formula's minus sign.
\end{remark}
